\begin{summarybox}
	\textbf{\textcolor{CSummary}{一句话核心}}
	\textbf{过双曲线外一点 $P$ 的两切线切点弦方程，只需将双曲线标准方程中的 $x^2$ 换为 $x_0 x$，$y^2$ 换为 $y_0 y$，即得。}

	\medskip
	\textbf{\textcolor{CSummary}{使用条件}}
	\begin{itemize}[leftmargin=*, itemsep=0.5em, topsep=0.4em]
		\item \textbf{条件1（点P外部）：} 点 $P(x_0,y_0)$ 必须在双曲线外部且能引两条实切线。
		\item \textbf{条件2（方程标准型）：} 双曲线方程为标准形式 $\frac{x^2}{a^2} - \frac{y^2}{b^2}=1$（焦点在 $x$ 轴）或 $\frac{y^2}{a^2} - \frac{x^2}{b^2}=1$（焦点在 $y$ 轴）。
	\end{itemize}

	\medskip
	\textbf{\textcolor{CSummary}{关键提醒}}
	\begin{itemize}[leftmargin=*, itemsep=0.5em, topsep=0.4em]
		\item \textbf{易错点（焦点匹配）：} 根据焦点位置正确选择公式，$x$ 轴用 $\frac{x_0 x}{a^2} - \frac{y_0 y}{b^2}=1$，$y$ 轴用 $\frac{y_0 y}{a^2} - \frac{x_0 x}{b^2}=1$。
		\item \textbf{检查项（符号减号）：} 方程中间为减号，不要误写为加号。
	\end{itemize}

\textbf{两有限实切点的存在条件。}
记 $q=x_0^2/a^2-y_0^2/b^2$，则恰有两个不同的有限实切点当且仅当
$q<1$ 且 $q\ne0$。$q=1$ 是一个切点，$q>1$ 无实切点；
$q=0$ 的非中心点在渐近线上，只有一个有限切点，另一方向为渐近线；中心点没有实切线。
纵轴形式使用 $q=y_0^2/a^2-x_0^2/b^2$。
\end{summarybox}
